% !TEX program = xelatex
\documentclass[11pt,a4paper]{article}

\usepackage[UTF8]{ctex}
\usepackage[margin=2.2cm]{geometry}
\usepackage{amsmath, amssymb, amsthm}
\usepackage{listings}
\usepackage{xcolor}
\usepackage{tikz}
\usepackage{booktabs}
\usepackage{multirow}
\usepackage{hyperref}
\usetikzlibrary{positioning, arrows.meta, shapes.geometric, fit, calc, decorations.pathreplacing}

\hypersetup{
  colorlinks=true,
  linkcolor=blue!60!black,
  urlcolor=teal!70!black,
}

\lstdefinelanguage{Rust}{
  morekeywords={fn,let,mut,pub,struct,impl,use,mod,self,Self,Option,Some,None,return,if,else,for,while,loop,match,break,continue,as,where,trait,type,enum,move,ref,const,static,unsafe,extern,crate,super,dyn,async,await,box,in},
  sensitive=true,
  morecomment=[l]{//},
  morecomment=[s]{/*}{*/},
  morestring=[b]",
}

\lstset{
  language=Rust,
  basicstyle=\ttfamily\footnotesize,
  keywordstyle=\color{blue!60!black}\bfseries,
  commentstyle=\color{green!40!black}\itshape,
  stringstyle=\color{red!60!black},
  numbers=left,
  numberstyle=\tiny\color{gray},
  numbersep=8pt,
  frame=single,
  framerule=0.4pt,
  rulecolor=\color{gray!40},
  breaklines=true,
  breakatwhitespace=true,
  showstringspaces=false,
  tabsize=4,
  columns=fullflexible,
  keepspaces=true,
}

\title{\textbf{halfedge\,: wgpu 几何导出模块设计文档} \\ \large \texttt{src/export.rs}}
\author{halfedge 库}
\date{2026-07-04}

\begin{document}
\maketitle
\tableofcontents

\section{模块定位}

\texttt{export.rs} 提供单一函数 \texttt{mesh\_to\_vertex\_index\_buffers}：
将半边网格导出为 \texttt{Vec<[f32; 3]>} 顶点缓冲与 \texttt{Vec<u32>} 索引缓冲，
可直接传入 \texttt{wgpu} 创建 \texttt{Buffer}。

\begin{center}
\renewcommand{\arraystretch}{1.18}
\begin{tabular}{@{}ll@{}}
  \toprule
  \textbf{API} & \texttt{mesh\_to\_vertex\_index\_buffers(\&MeshStorage) -> (Vec<[f32; 3]>, Vec<u32>)} \\
  \midrule
  顶点缓冲 & 每个顶点 3 个 \texttt{f32}（位置 x/y/z） \\
  索引缓冲 & 每个三角面 3 个 \texttt{u32}（CCW 朝向） \\
  索引类型 & \texttt{u32}（wgpu 默认 \texttt{IndexFormat::Uint32}） \\
  顶点编号 & 按 \texttt{vertex\_ids()} 顺序 0-based 重新编号 \\
  \bottomrule
\end{tabular}
\end{center}

\section{设计目标}

\subsection{为何需要单独的导出模块？}

半边网格适合拓扑查询与编辑，但 \texttt{wgpu} 渲染需要扁平的顶点+索引缓冲：
\begin{itemize}
  \item 顶点缓冲：连续内存，每个顶点固定字节数（此处 12 字节 = 3 × f32）；
  \item 索引缓冲：用顶点索引引用三角形，避免重复存储共享顶点。
\end{itemize}

\texttt{export.rs} 负责「半边网格 $\to$ 扁平缓冲」的转换，是渲染管线的入口。

\subsection{为何不用 \texttt{Vec<f32>} 而用 \texttt{Vec<[f32; 3]>}？}

\texttt{[f32; 3]} 在内存中与 3 个连续 \texttt{f32} 完全等价（12 字节，4 字节对齐），
但语义更清晰：每个元素代表一个顶点位置。调用方可直接 \texttt{as\_ptr()} 转
\texttt{*const f32} 传给 \texttt{wgpu::Buffer}，也可方便地 \texttt{extend} 法向、
UV 等属性。

\section{导出算法}

\subsection{流程图}

\begin{center}
\resizebox{\textwidth}{!}{%
\begin{tikzpicture}[
  node distance=0.5cm,
  proc/.style={draw, rounded corners=2pt, minimum width=10cm, minimum height=0.7cm, align=center, font=\small},
  op/.style={-Stealth, thick},
  startend/.style={draw, rounded corners=10pt, minimum width=2.6cm, minimum height=0.7cm, font=\small, fill=gray!12}
]
  \node[startend] (s) {开始：\&MeshStorage};
  \node[proc, below=of s, fill=blue!8] (p1) {阶段 1：遍历 \texttt{vertex\_ids()}，构建 \texttt{HashMap<VertexId, u32>} 与 \texttt{Vec<[f32; 3]>}};
  \node[proc, below=of p1, fill=blue!8] (p2) {阶段 2：遍历 \texttt{face\_ids()}，每面取 3 个 tip 顶点，查映射输出 \texttt{u32} 索引};
  \node[proc, below=of p2, fill=yellow!15] (p3) {跳过：非三角面 / 顶点缺失};
  \node[proc, below=of p3, fill=green!12] (p4) {返回 \texttt{(vertices, indices)}};
  \node[startend, below=of p4] (e) {结束};
  \draw[op] (s) -- (p1);
  \draw[op] (p1) -- (p2);
  \draw[op] (p2) -- (p3);
  \draw[op] (p3) -- (p4);
  \draw[op] (p4) -- (e);
\end{tikzpicture}
}
\end{center}

\subsection{阶段 1：顶点收集与索引映射}

\begin{itemize}
  \item 遍历 \texttt{mesh.vertex\_ids()}（slotmap 顺序）；
  \item 对每个 \texttt{VertexId}，取其位置 \texttt{[f64; 3]}，截断为 \texttt{[f32; 3]}；
  \item 用 \texttt{HashMap<VertexId, u32>} 记录「VertexId $\to$ 0-based 索引」映射；
  \item 顶点缺失（已删除）跳过。
\end{itemize}

\subsection{阶段 2：索引生成}

\begin{itemize}
  \item 遍历 \texttt{mesh.face\_ids()}；
  \item 对每个面，用 \texttt{FaceHalfEdges} 迭代器取边界环半边；
  \item 取每条半边的 \texttt{vertex}（tip），查 \texttt{HashMap} 得 0-based 索引；
  \item 边界环长度 $\ne 3$ 跳过（非三角面）；
  \item 索引按 CCW 顺序输出（与原面 \texttt{next} 链一致）。
\end{itemize}

\section{精度与兼容性}

\subsection{f64 $\to$ f32 截断}

半边网格内部用 \texttt{f64} 存储顶点位置（精度高，适合几何计算）；
导出时截断为 \texttt{f32}（GPU 原生支持，性能优）。

\textbf{精度损失}：\texttt{f64} 有约 15--17 位有效数字，\texttt{f32} 约 6--9 位。
对大多数渲染场景（屏幕分辨率、视锥范围）足够；若需要高精度渲染
（如 CAD 模型远距离观察），可扩展为 \texttt{Vec<[f32; 6]>}（双精度拆分）或
使用 \texttt{wgpu::VertexAttribute} 的 \texttt{Float64} 格式。

\subsection{wgpu 集成示例}

\begin{lstlisting}
use wgpu::util::DeviceExt;

let (vertices, indices) = mesh_to_vertex_index_buffers(&mesh);

let vertex_buffer = device.create_buffer_init(
    &wgpu::util::BufferInitDescriptor {
        label: Some("Vertex Buffer"),
        contents: bytemuck::cast_slice(&vertices),
        usage: wgpu::BufferUsages::VERTEX,
    },
);

let index_buffer = device.create_buffer_init(
    &wgpu::util::BufferInitDescriptor {
        label: Some("Index Buffer"),
        contents: bytemuck::cast_slice(&indices),
        usage: wgpu::BufferUsages::INDEX,
    },
);

// 渲染时：
// render_pass.set_vertex_buffer(0, vertex_buffer.slice(..));
// render_pass.set_index_buffer(index_buffer.slice(..), wgpu::IndexFormat::Uint32);
// render_pass.draw_indexed(0..indices.len() as u32, 0, 0..1);
\end{lstlisting}

\section{复杂度}

$O(V + F)$：每个顶点与每个面常数时间。

\begin{center}
\renewcommand{\arraystretch}{1.18}
\begin{tabular}{@{}lll@{}}
  \toprule
  \textbf{阶段} & \textbf{时间} & \textbf{空间} \\
  \midrule
  顶点收集 & $O(V)$ & $O(V)$（\texttt{Vec} + \texttt{HashMap}） \\
  索引生成 & $O(F)$ & $O(F)$（\texttt{Vec<u32>}） \\
  总计     & $O(V + F)$ & $O(V + F)$ \\
  \bottomrule
\end{tabular}
\end{center}

\section{退化守卫}

\begin{itemize}
  \item 顶点缺失（\texttt{get\_vertex} 返回 \texttt{None}）：跳过，不输出；
  \item 面边界环长度 $\ne 3$：跳过该面，不输出索引；
  \item 半边 \texttt{vertex} 在 \texttt{HashMap} 中找不到（理论不应发生，
        因阶段 1 已收集所有顶点）：\texttt{filter\_map} 静默跳过。
\end{itemize}

\section{测试覆盖}

\begin{center}
\renewcommand{\arraystretch}{1.18}
\begin{tabular}{@{}ll@{}}
  \toprule
  \textbf{测试名} & \textbf{验证点} \\
  \midrule
  \texttt{export\_basic\_quad}             & 4 顶点 6 索引，位置正确 \\
  \texttt{export\_tetrahedron}             & 4 顶点 12 索引（4 面） \\
  \texttt{export\_preserves\_ccw\_winding} & 导出索引 CCW 朝向，法向 $+z$ \\
  \texttt{export\_empty\_mesh}             & 空网格输出空缓冲 \\
  \texttt{export\_roundtrip\_validation}   & 导出 $\to$ 重建 $\to$ 通过校验 \\
  \bottomrule
\end{tabular}
\end{center}

\end{document}
